% !TEX root = ../../tree_notes.tex
% =====================================================================
%  Chapter 6 opener -- "Chapter at a Glance", reading paths + ladder
%  staircase, the running example T seen through FOUR data shadows
%  (edge list, parent array, degree sequence, queries), roadmap,
%  recognition guide, constraint table, notation.  NO numbered boxes here.
% =====================================================================

\begin{chaptertheorybox}{Chapter at a Glance}
Chapters~4 and~5 pushed \emph{forward}: here is a tree, store it, then lay it on a line.
This chapter pulls \emph{backward}. A contest hands you \term{raw data} -- a list of edges, a
\term{parent array}, a \term{degree sequence}, an \term{encoding}, even \emph{answers to
questions you ask} -- and the real task is to \term{rebuild the tree} and then \emph{prove the
rebuild is legitimate}: connected, acyclic, exactly $n-1$ edges, a single root, degrees summing
to $2(n-1)$. That reverse move -- \emph{identify the structure hiding inside the data} -- is the
skill every later chapter silently assumes the moment it reads \code{edges[i]=[u,v]}. The three
decisions this chapter forces on you:
\begin{itemize}[nosep]
  \item \textbf{Given pairs, is there a tree, and which one?} An edge list or parent array
        determines a tree \emph{only} if it passes the \term{tree certificate}: $n$ nodes,
        $n-1$ edges, connected \emph{and} acyclic (any two of the three force the third). Root
        it by finding the node with no parent. \hfill \lv{2}--\lv{4}
  \item \textbf{Given degrees or a code, can a tree realise them?} A list of positive integers
        is some tree's degree sequence \emph{iff} it sums to $2(n-1)$; the \term{Pr\"ufer code}
        turns every labelled tree into a length-$(n-2)$ sequence, so a degree array even tells
        you \emph{how many} trees exist. \hfill \lv{3}--\lv{5}
  \item \textbf{Given only yes/no and numbers from an oracle, can you still find it?} Ask a
        \emph{distance} or \emph{ancestor} query, get one integer or bit back, and adapt. Each
        honest answer halves the candidates -- reconstructing a hidden tree is binary search on
        a shape. \hfill \lv{2}--\lv{4}
\end{itemize}

\textbf{You will learn to}
\begin{itemize}[nosep]
  \item rebuild an adjacency structure from an edge list or a parent array and locate the root
        in one pass; \hfill\lv{2}--\lv{3}
  \item certify that a graph \emph{is} a tree using the $n-1$-edges / connected / acyclic
        triangle, with a union-find pre-look; \hfill\lv{3}--\lv{4}
  \item decide whether a degree sequence is realisable and construct a tree that realises it;
        \hfill\lv{3}--\lv{4}
  \item read the Pr\"ufer code as the bridge between trees, degree arrays and counting
        (full bijection in \chref{c32}); \hfill\lv{4}--\lv{5}
  \item repair an over-full edge set to a tree (delete the redundant edge) and recover a hidden
        structure from adaptive queries. \hfill\lv{2}--\lv{5}
\end{itemize}

\textbf{Prerequisites.} Representation and rooting from Chapter~4 (\secref{c04-rooting},
\coderef{c04-root}) -- you already know how to pour edges into an \code{adj} list and DFS with a
parent guard; this chapter adds the \emph{validity} and \emph{reconstruction} decisions on top.
The timestamp/distance machinery of Chapter~5 (\chref{c24}) is \emph{used} as an oracle in
\secref{c25-query}, never re-derived.
\end{chaptertheorybox}

% ---------------- Station 1b: reading paths + ladder position ----------------
\begin{chaptertheorybox}{How to Read This Chapter (two paths up the ladder)}
Level badges $\lv1$--$\lv6$ as before: $\lv2$ is ``read the data, rebuild the array'', $\lv3$ is
the solid interview medium, $\lv5$ is an LC Hard reconstruction with a case split.

\textbf{Beginner path ($\lv2$--$\lv3$).} Read \secref{c25-parent} fully -- it turns an edge list
or parent array into a rooted tree and states the tree certificate. Then \secref{c25-identify}
up through the union-find ``is it a tree?'' test. Take every $\lv2$--$\lv3$ solved problem. You
may \emph{skip on a first pass}: the degree/Pr\"ufer material of \secref{c25-degree}, the
interactive protocol of \secref{c25-query}, every $\lv4{+}$ variation, the hardening notes and
the Summit. At the end of this path you can rebuild any tree from an edge list and certify it in
$O(n)$.

\textbf{Full path.} Everything, in order -- it ends with the scaled Summit (LC Hard 1719
\emph{Number of Ways to Reconstruct a Tree}: read an ancestor-pair table, sort nodes by degree,
greedily attach each node to the smallest feasible parent, and decide \emph{unique / many /
impossible}) and the reduction gallery showing that ``rebuild from a parent array'', ``realise a
degree sequence'', ``delete the redundant edge'' and ``find the node by distance queries'' are the
\emph{same} certificate checked against four different data sources. The heavier theory -- the
Disjoint-Set-Union structure itself and the Pr\"ufer/Cayley counting proofs -- lives in
\chref{c28} and \chref{c32}.
\end{chaptertheorybox}

\begin{figure}[H]
\centering
\begin{tikzpicture}[diagram, x=9mm, y=7mm]
  \foreach \st/\lab in {1/lv1,2/lv2,3/lv3,4/lv4,5/lv5,6/lv6} {
    \fill[\lab!22, draw=\lab, line width=0.8pt] (\st-1,0) rectangle (\st,\st*0.52);
    \node[font=\scriptsize\bfseries, text=\lab] at (\st-0.5,\st*0.52-0.42) {L\st};}
  \node[annot, anchor=south west, font=\scriptsize]
    at (0.05,0.62) {6.1 read data, rebuild adj};
  \node[annot, anchor=south west, font=\scriptsize]
    at (1.05,1.14) {6.1/6.3 tree certificate};
  \node[annot, anchor=south west, font=\scriptsize]
    at (2.05,1.66) {6.3 union-find pre-look};
  \node[annot, anchor=south west, font=\scriptsize]
    at (3.05,2.18) {6.2 degree sequence, Pr\"ufer teaser};
  \node[annot, anchor=south west, font=\scriptsize, text=lv5]
    at (4.05,2.70) {6.2 Pr\"ufer counting, 6.4 adaptive queries};
  \node[annot, anchor=south west, font=\scriptsize, text=lv5]
    at (5.05,3.22) {scaled summit (1719)};
\end{tikzpicture}
\caption*{Chapter~6 on the six-step ladder: rebuilding and certifying carry you from raw pairs to
``this is a tree and here is its root'' ($\lv3$); the degree/Pr\"ufer and interactive topics and the
scaled Summit (1719) reach $\lv5$. Beginner path $=$ steps 2--3.}
\end{figure}

% ---------------- The running example ----------------
\subsection*{The Running Example: One Hidden Tree, Four Data Shadows}

This chapter has \emph{one} running example: the general tree $T$ below, drawn once and then
\emph{never shown again} -- instead every topic receives it as a different pile of \emph{data}.
The point is that four unrelated-looking inputs (an edge list, a parent array, a degree sequence,
a bag of distance answers) all describe \emph{this exact tree}, and reconstructing it is the whole
game. Rooted at node $1$, children visited in ascending order:
\[
  \text{edges } \{1\text{--}2,\,1\text{--}3,\,2\text{--}4,\,2\text{--}5,\,3\text{--}6,\,3\text{--}7,\,7\text{--}8\}.
\]

\begin{figure}[H]
\centering
\begin{tikzpicture}[diagram]
  \node[troot] (n1) at (0,0) {$1$};
  \node[tnode] (n2) at (-3,-1.4) {$2$};
  \node[tnode] (n3) at (3,-1.4) {$3$};
  \node[tleaf] (n4) at (-4,-2.8) {$4$};
  \node[tleaf] (n5) at (-2,-2.8) {$5$};
  \node[tleaf] (n6) at (2,-2.8) {$6$};
  \node[tnode] (n7) at (4,-2.8) {$7$};
  \node[tleaf] (n8) at (4,-4.2) {$8$};
  \draw[tedgehl] (n1) -- (n2); \draw[tedgehl] (n1) -- (n3);
  \draw[tedgehl] (n2) -- (n4); \draw[tedgehl] (n2) -- (n5);
  \draw[tedgehl] (n3) -- (n6); \draw[tedgehl] (n3) -- (n7);
  \draw[tedgehl] (n7) -- (n8);
  \node[annot, align=left, font=\scriptsize] at (-6.6,-1.4)
    {$n=8$, $7=n-1$ edges\\ root $=1$\\ leaves $=\{4,5,6,8\}$};
  \node[annot, align=left, font=\scriptsize] at (6.4,-2.0)
    {diameter $4\!\to\!8$\\ is $5$ edges\\ height $=3$};
\end{tikzpicture}
\caption{Running example $T$ (rooted at $1$). Four different data sources -- the edge list above,
the parent array, the degree sequence and a set of distance answers -- each alone is enough to
rebuild this tree, and no two of them look alike. That is the chapter in one picture.}\label{fig:c25-running}
\end{figure}

\medskip
\noindent Keep these four ``shadows'' of $T$ in mind -- they are re-derived, never merely asserted,
in \secref{c25-parent} and \secref{c25-degree}:
\begin{center}
\renewcommand{\arraystretch}{1.25}
\setlength{\tabcolsep}{5pt}
\begin{tabular}{@{}p{0.24\linewidth} p{0.70\linewidth}@{}}
\toprule
\textbf{Data shadow} & \textbf{Value for $T$ (rooted at $1$, nodes $1..8$)} \\
\midrule
edge list (undirected)   & $[1,2],[1,3],[2,4],[2,5],[3,6],[3,7],[7,8]$ -- $n-1=7$ pairs \\
parent array $\parent[\cdot]$ & $[\,0,1,1,2,2,3,3,7\,]$ (index $i$ is node $i$; $\parent[1]=0$ marks the root) \\
degree sequence $\degop[\cdot]$ & $[\,2,3,3,1,1,1,2,1\,]$, sum $=14=2(n-1)$ \\
Pr\"ufer code (teaser)    & $(2,2,1,3,3,7)$ -- length $n-2=6$; node $v$ appears $\degop(v)-1$ times \\
distance answers (6.4)   & $\dist(4,8)=5$, $\dist(5,6)=4$, $\dist(1,8)=3$, \ldots\\
\bottomrule
\end{tabular}
\end{center}

% ---------------- Roadmap diagram ----------------
\begin{figure}[H]
\centering
\begin{tikzpicture}[diagram, node distance=5mm,
    top/.style={boxnode, fill=primaryblue!10, draw=primaryblue, text width=25mm,
                minimum height=11mm, font=\small},
    res/.style={annot, text width=25mm, anchor=north}]
  \node[top] (t1) {\textbf{6.1} Parent / Edge};
  \node[top, right=of t1] (t2) {\textbf{6.2} Degrees \& Codes};
  \node[top, right=of t2] (t3) {\textbf{6.3} Is It a Tree?};
  \node[top, right=of t3] (t4) {\textbf{6.4} By Queries};
  \foreach \a/\b in {t1/t2, t2/t3, t3/t4} \draw[flow] (\a) -- (\b);
  \node[res] at ([yshift=-2mm]t1.south) {rebuild adj; find root; certify $n{-}1$ edges, connected, acyclic};
  \node[res] at ([yshift=-2mm]t2.south) {realisable degree sequence?; greedy build; Pr\"ufer teaser};
  \node[res] at ([yshift=-2mm]t3.south) {union-find pre-look; tree, forest or cycle; drop one edge};
  \node[res] at ([yshift=-2mm]t4.south) {query model; one answer halves it; adaptive search};
\end{tikzpicture}
\caption*{The four topics, left to right, are the \emph{same certificate} tested against four data
sources. 6.1 rebuilds from \emph{pairs}; 6.2 from \emph{degrees/codes}; 6.3 \emph{certifies} (and
repairs) a suspicious edge set; 6.4 recovers the shape from \emph{answers}. Each is construction
from data seen one way.}
\end{figure}

% ---------------- Topic table ----------------
\begin{center}
\renewcommand{\arraystretch}{1.25}
\setlength{\tabcolsep}{4pt}
\begin{tabular}{@{}p{0.24\linewidth} p{0.36\linewidth} c p{0.23\linewidth}@{}}
\toprule
\textbf{Topic} & \textbf{Key results} & \textbf{Levels} & \textbf{Solved problems} \\
\midrule
\secref{c25-parent}\ \ Parent Array / Edge List
  & rebuild adjacency; root-finding; the tree certificate; repair an over-edge set & $\lv2$--$\lv4$
  & \probref{c25-rc2} \\
\addlinespace
\secref{c25-degree}\ \ Degrees, Sequences \& Codes
  & degree sequence realisable $\iff$ sum $2(n{-}1)$; greedy build; Pr\"ufer teaser & $\lv3$--$\lv5$
  & \probref{c25-quad} \\
\addlinespace
\secref{c25-identify}\ \ Is It a Tree? (Union-Find)
  & connected + acyclic + $n{-}1$ edges triangle; union-find test; tree/forest/cyclic & $\lv2$--$\lv3$
  & \probref{c25-rc1} \\
\addlinespace
\secref{c25-query}\ \ Recover by Queries
  & interactive model; what one query reveals; adaptive halving for a hidden node & $\lv2$--$\lv4$
  & \probref{c25-guess} \\
\bottomrule
\end{tabular}
\end{center}

% ---------------- Problem recognition guide ----------------
\subsection*{Problem Recognition Guide --- When You See \emph{this}, Think \emph{that}}

Construction problems announce themselves by \emph{withholding the drawing}: you get pairs,
counts, an array or an oracle instead of pointers. Match the phrasing; every row points at the
section that builds the tool.

\begin{center}
\renewcommand{\arraystretch}{1.3}
\setlength{\tabcolsep}{4pt}
\begin{tabular}{@{}p{0.42\linewidth} p{0.35\linewidth} l@{}}
\toprule
\textbf{When the problem says\ldots} & \textbf{\textcolor{deepblue}{\textbf{Think\ldots}}} & \textbf{Section} \\
\midrule
``\code{edges[i]=[u,v]}'', ``\code{parent[i]} gives the parent of $i$'', a forest given as pairs
  & rebuild \code{adj}, find the root, verify the certificate & \secref{c25-parent} \\
\addlinespace
``is it a valid tree?'', ``one extra edge added -- remove it'', ``which edge is redundant''
  & the $n{-}1$/connected/acyclic test; union-find pre-look & \secref{c25-identify} \\
\addlinespace
``given the degrees of every node, build a tree'', ``how many trees match this degree array''
  & degree sequence sums to $2(n-1)$; Pr\"ufer code counts them & \secref{c25-degree} \\
\addlinespace
``an encoding / grid / nested string -- construct the tree''
  & recursive build-by-rule (quad tree); a stack parse & \secref{c25-degree} \\
\addlinespace
``you may ask \emph{queries}; find the hidden node/edge in $k$ questions''
  & each answer is a bit -- adapt, halve, descend & \secref{c25-query} \\
\addlinespace
``reconstruct so that every listed ancestor-pair holds; is it unique?''
  & sort by degree, greedy parent, transitivity check (Summit) & \probref{c25-summit} \\
\bottomrule
\end{tabular}
\end{center}

\begin{center}
\renewcommand{\arraystretch}{1.25}
\setlength{\tabcolsep}{4pt}
\begin{tabular}{@{}p{0.27\linewidth} p{0.45\linewidth} p{0.20\linewidth}@{}}
\toprule
\textbf{Constraint you see} & \textbf{What it forces} & \textbf{Chapter-6 tool} \\
\midrule
$n-1$ edges promised, graph may be corrupt & do not trust it -- run the certificate/union-find & \secref{c25-identify} \\
a tree \emph{reconstructed} then \emph{queried} many times & rebuild once from the array, answer later & \secref{c25-parent} \\
``how many / which structure fits this degree list'' & Pr\"ufer bijection: degrees pin the count & \secref{c25-degree} (\chref{c32}) \\
interactive judge, tight query budget & information-per-query drives the strategy & \secref{c25-query} \\
$N$ up to $2\cdot10^{5}$ edge list & rebuild \emph{iteratively}; $\bigO{n\alpha}$ union-find & \secref{c25-identify} (\chref{c28}) \\
\bottomrule
\end{tabular}
\end{center}

\begin{notebox}[Not this chapter?]
If the tree is \emph{already given} as pointers or an \code{adj} list and you only need to walk or
store it, that is Chapter~4 (\chref{c04}). If you must rebuild a \emph{binary} tree from two of its
\emph{traversal orders} (inorder+preorder, inorder+postorder) or serialise/parse it, that is
Chapter~3 (\chref{c03}) -- construction from \emph{traversals}, a different data source than the
\emph{arbitrary} data here. If the real question is once the tree stands -- subtree ranges,
distances, LCA -- that is Chapter~5 (\chref{c24}) and Chapter~16 (\chref{c11}). If you need the
Disjoint-Set-Union structure itself (path compression, rank), that is Chapter~17
(\chref{c28}); we borrow only its \emph{pre-look} as a tree test. If you want the full Pr\"ufer
bijection and Cayley's $n^{n-2}$ \emph{proof}, that is Chapter~30 (\chref{c32}). This chapter owns
\emph{reading a tree out of whatever data you were handed, and proving you read it right}.
\end{notebox}

% ---------------- Notation ----------------
\subsection*{Notation Used in This Chapter}

\begin{center}
\renewcommand{\arraystretch}{1.2}
\setlength{\tabcolsep}{4pt}
\begin{tabular}{@{}p{0.20\linewidth} p{0.27\linewidth}@{\quad}p{0.17\linewidth} p{0.27\linewidth}@{}}
\toprule
\textbf{Symbol} & \textbf{Meaning} & \textbf{Symbol} & \textbf{Meaning} \\
\midrule
$T$                 & the running general tree (\figref{c25-running}) & $\parent(v)$ & $v$'s parent; $0$/$-1$ marks the root \\
$\degop(v)$         & number of neighbours of $v$ (degree)           & $\mathcal{E}$ & edge set, $|\mathcal{E}|=n-1$ for a tree \\
$\mathcal{P}$       & the Pr\"ufer code, length $n-2$                & $\mathcal{C}$ & number of connected components \\
$\dist(u,v)$        & edges on the $u$--$v$ path (from \chref{c24})  & $\mathcal{R}$ & a query/answer oracle for a hidden tree \\
$\mathrm{find},\mathrm{uni}$ & union-find ops (\chref{c28})     & $\lv{3}$       & difficulty badge (runs $\lv{1}$--$\lv{6}$) \\
\bottomrule
\end{tabular}
\end{center}

As in earlier chapters, $\dist$ and $\depth$ count \textbf{edges}, general trees are stored as an
adjacency list \code{adj} (\code{vector} of \code{vector<int>}) with a parent guard, and every
listing is complete $\cpp$ (standard header included). Node numbering follows the problem: many
tree-input problems are \textbf{$1$-indexed} (LeetCode 684/685 label nodes $1..n$), which is why
the running example $T$ is $1$-indexed here; \code{adj} arrays are $0$- or $1$-sized to match.
